\lecture{25}{Jun 5.}{}

Recall the primal program above is 
\[
    \begin{aligned}
        \max~ & \vv{c}^T \vv{x}  \\
            & A \vv{x} \le \vv{0}  \\
            & \vv{x} \ge \vv{0},
    \end{aligned}
\]
where the dual LP is 
\[
    \begin{aligned}
        \min~ & \vv{b}^T \vv{y}  \\
            & \vv{A}^T \vv{y} \ge \vv{e}  \\
            & \vv{y} \ge \vv{0}.
    \end{aligned}
\]

\begin{remark}
    If some \(x_i\) can be negative, then the dual program above is not sufficient. 
\end{remark}

\begin{lemma}[Weak Duality]
    If \(\vv{x}\) is feasible for the primal and \(\vv{y}\) is feasible for the dual, then 
    \[
        \vv{c}^T \vv{x} \le \vv{b}^T \vv{y}.
    \]  
\end{lemma}

\begin{theorem}[Strong Duality]
    Given a primal LP and its dual, there are four possibilitiws:
    \begin{enumerate}[(i)]
        \item both are infeasible. 
        \item primal infeasible, dual unbounded. 
        \item primla unbounded, dual infeasible. 
        \item both are feasible and bounded, and if \(\vv{x}\) is optimal for the primal and \(y\) is optimal for the dual, then 
        \[
            \vv{c}^T \vv{x} = \vv{b}^T \vv{y}.
        \]
    \end{enumerate}
\end{theorem}

\begin{theorem}[Standard form]
    Every linear program can be written in the form 
    \[
        \begin{aligned}
            \max~ & \vv{c}^T \vv{x}  \\
                & A \vv{x} = \vv{b}  \\
                & \vv{x} \ge \vv{0}.
        \end{aligned}
    \]
\end{theorem}
\begin{proof}
    If we have 
    \[
        a_{i1} x_1 + a_{i2} x_2 + \dots + a_{in} x_n \le b_i,
    \]
    we turn this into equation by intoducing a new slack variable 
    \[
        a_{i1} x_1 + a_{i2} x_2 + \dots + a_{in} x_n + z_i = b_i, \quad z_i \ge 0.
    \]

    If we have \(x_i \in \mathbb{R}\) instead of \(x_i \ge 0\). We replace \(x_i\) with \(x_i^+ - x_i^-\), where \(x_i^+ , x_i^- \ge 0\). (here \(x_i^+ = \max (x_i, 0), x_i^- = \max (-x_i, 0)\))    
\end{proof}

\section{Simplex method}

LP with \(n\) variables, \(m\) equality constraints, where \(n \ge m\), has
\begin{itemize}
    \item feasible set is a polyhedron 
    \item optimal value attained at a vertex
    \item Simplex algorithm: walks along the vertices of this polyhedron. Keep moving to a neighboring vertex of higher objective value.
\end{itemize}   

Simplex solves primal and dual simultaneously maubtaub \(\vv{x}\) for the primal and \(\vv{y}\) for the dual, where
\begin{itemize}
    \item \(\vv{c}^T \vv{x} = \vv{b}^T \vv{y}\)
    \item \(\vv{x}\) is feasible for primal.  
\end{itemize} 

Thus, we know we are optimal as soon as \(\vv{y}\) is feasible for the dual. 

\begin{question}
    How do we move between vertices?
\end{question}
\begin{answer}
    Called "pivoting" in the algorithm, which specified by the "pivot rule".
\end{answer}

For every common pivot rule, there are examples where the Simplex algorithm takes exponential time. 

\begin{question}[Open problem]
    Is there a poly-time pivot rule?
\end{question}

%fig

\paragraph{Interior Point Methods(Ellipsoid Algorithm)}
Poly-time algorithm for solving LPs: move through the interior of the feasible set.

\section{Integer linear programming}
Linear programs with integer variables, which is very useful for combinatorial problrms (by using 0/1 indicator variables), and ILP is NP complete to solve.

\paragraph{LP relaxion}
Take an ILP, and forget that the variables should be integers, then LP can be solved efficiently, and we can bound on the original ILP problem.

\begin{eg}[Fractional chromatic number]
    Original problem is an ILP for computing \(\chi (G)\): 
    \[
        \begin{aligned}
            \min~ & \sum_{I \in I(G)} \chi _I   \\
                & \forall v \in V(G), \sum_{I \ni v} \chi _I \ge 1   \\
                & \chi _I \in \left\{ 0, 1 \right\} \quad \forall I \in I(G). 
        \end{aligned}
    \]
    Then we can relax it the last constraint to \(\chi _I \ge 0\) and \(\chi (G) \ge \chi _f (G)\), where \(\chi _f(G)\) is the fractional chromatic number.   
\end{eg}

\begin{eg}
    \(\chi (C_5) = 3\), and \(\chi _I = \frac{1}{2}\) for all \(I, \vert I \vert = 2, \) and independent, is a relaxed estimation to \(\chi (C_5)\), which gives \(\sum_{I} \chi _I = \frac{5}{2} < 3 \), an lower bound for \(\chi (C_5)\).    
\end{eg}